\documentclass[conference]{IEEEtran}
\usepackage[utf8]{inputenc}
\usepackage[T1]{fontenc}
\usepackage{cite}
\usepackage{amsmath,amssymb,amsfonts}
\usepackage{algorithmic}
\usepackage{graphicx}
\usepackage{textcomp}
\usepackage{xcolor}
\usepackage{booktabs}
\usepackage{hyperref}
\usepackage{float}
\usepackage{url}
\usepackage{balance}
\usepackage{flushend}
\graphicspath{{./images/}{./}} 

\setlength{\textfloatsep}{4pt plus 1pt minus 1pt}
\setlength{\intextsep}{4pt plus 1pt minus 1pt}
\setlength{\floatsep}{4pt plus 1pt minus 1pt}
\setlength{\abovedisplayskip}{3pt}
\setlength{\belowdisplayskip}{3pt}
\setlength{\abovedisplayshortskip}{0pt}
\setlength{\belowdisplayshortskip}{3pt}

\linespread{0.98}

\usepackage{enumitem}
\setlist{nosep, leftmargin=*}

\widowpenalty=10000
\clubpenalty=10000


\begin{document}

\title{Gera\c{c}\~ao Sint\'etica Guiada por Heran\c{c}a Ativa para Detec\c{c}\~ao de Textos Fraudulentos}

\author{
    \IEEEauthorblockN{Eduardo Alexandre de Amorim\textsuperscript{1}, Cleber Zanchettin\textsuperscript{1}}
    \IEEEauthorblockA{\textsuperscript{1}Centro de Inform\'atica, Universidade Federal de Pernambuco, Recife, PE, Brazil}
    \IEEEauthorblockA{\{eaa5, cz\}@cin.ufpe.br}
}

\maketitle

% --- AVISO DE COPYRIGHT PARA O ARXIV ---
\def\thefootnote{}\footnotetext{\copyright~2026 IEEE. Personal use of this material is permitted. Permission from IEEE must be obtained for all other uses, in any current or future media, including reprinting/republishing this material for advertising or promotional purposes, creating new collective works, for resale or redistribution to servers or lists, or reuse of any copyrighted component of this work in other works.}
% ---------------------------------------


\begin{abstract}
A escassez e o desbalanceamento de dados reais em ciberseguran\c{c}a dificultam o treinamento de classificadores robustos contra textos fraudulentos. A gera\c{c}\~ao de dados sint\'eticos via modelos de linguagem \'e uma alternativa, mas processos gerativos est\'aticos podem causar colapso de distribui\c{c}\~ao e baixa cobertura de casos raros. Este trabalho investiga a \emph{Heran\c{c}a Ativa}: um loop gerativo c\'iclico em que um avaliador adversarial prioriza amostras sint\'eticas com base no erro, na incerteza e na diversidade de um classificador. Em experimentos com o corpus Fake.br em cen\'ario de escassez (10\% da classe fraudulenta), o m\'etodo alcan\c{c}ou F1 de $0.924$, superior \`a gera\c{c}\~ao est\'atica ($0.916$) e ao treino sem dados sint\'eticos ($0.887$), com a menor vari\^ancia entre sementes (desvio padr\~ao $0.007$). Entretanto, a hip\'otese de efici\^encia computacional n\~ao foi confirmada: o m\'etodo consome mais tokens que a gera\c{c}\~ao est\'atica para atingir o mesmo F1.
\end{abstract}

\begin{IEEEkeywords}
synthetic data generation, large language models, fake news detection, data-centric AI, active inheritance
\end{IEEEkeywords}

\section*{Reproducibility}
C\'odigo do projeto: \url{https://github.com/Edu-p/active-inheritance-fraud}. Baseline: \url{https://github.com/szanara/profile2gen}. Dados: Fake.br-Corpus ($7{,}200$ not\'icias, $3{,}600$ falsas e $3{,}600$ verdadeiras). Cen\'ario de escassez: treino com $10\%$ da classe falsa ($288$ falsas, $2{,}880$ verdadeiras) e holdout balanceado intacto ($720$/$720$). Classificador: BERTimbau (\texttt{neuralmind/bert-base-portuguese-cased}). Embeddings: MiniLM multil\'ingue (\texttt{paraphrase-multilingual-MiniLM-L12-v2}). Gerador: \texttt{deepseek-chat}. Sementes: $42$, $43$, $44$. GPU: NVIDIA RTX 4060 (8\,GB).

\section{Introdu\c{c}\~ao}
A detec\c{c}\~ao de textos fraudulentos (not\'icias falsas, phishing, an\'uncios enganosos) sofre com a escassez e o desbalanceamento de dados reais~\cite{souza2025}. Embora a gera\c{c}\~ao de dados sint\'eticos via modelos de linguagem seja uma alternativa, processos est\'aticos frequentemente resultam em colapso do modelo e baixa cobertura de casos raros~\cite{shumailov2024}. Este trabalho adapta o framework \emph{Profile2Gen}~\cite{souza2025} --- originalmente proposto para dados tabulares de regress\~ao --- ao dom\'inio textual de classifica\c{c}\~ao, introduzindo um \emph{avaliador adversarial} no loop gerativo. A diferen\c{c}a conceitual central \'e uma invers\~ao do baseline: enquanto o Profile2Gen \emph{remove} amostras dif\'iceis no p\'os-processamento, a Heran\c{c}a Ativa \emph{herda} as amostras dif\'iceis/incertas como ``treino deliberado'' para o classificador.

As contribui\c{c}\~oes s\~ao: (i) um loop gerativo c\'iclico guiado por erro, incerteza e diversidade; (ii) uma avalia\c{c}\~ao no corpus Fake.br em cen\'ario de escassez; e (iii) uma an\'alise das dimens\~oes de robustez, diversidade e custo computacional.

\section{Solu\c{c}\~ao Proposta}
O pipeline (Algoritmo~\ref{alg:ai}) opera em ciclos. A cada ciclo $t$, o LLM gera um lote $S_t$ condicionado ao pool herdado $H$ (few-shot). Um classificador (\emph{learner}) produz $p(y|x)$ e o erro por amostra; um codificador produz embeddings para a diversidade. O avaliador adversarial atribui a cada amostra um escore multiobjetivo:

\begin{equation}
score(x) = w_1 \cdot err(x) + w_2 \cdot unc(x) + w_3 \cdot div(x),
\end{equation}

onde $err$ \'e o log-loss do learner, $unc$ \'e a entropia normalizada de $p(y|x)$, e $div$ \'e a dist\^ancia de cosseno m\'axima ao pool herdado (diversidade). As $k$ amostras de maior escore (fronteira de decis\~ao) s\~ao herdadas para o pr\'oximo ciclo e injetadas no treino. Usamos $w_1{=}w_2{=}0.4$ e $w_3{=}0.2$.

\begin{figure}[t]
\caption{Heran\c{c}a Ativa (loop c\'iclico).}
\label{alg:ai}
\begin{algorithmic}[1]
\STATE Treina learner $f$ em Data 0 (treino escasso)
\STATE Pool herdado $H \leftarrow \emptyset$
\FOR{$t = 0$ \TO $T-1$}
  \STATE Gera lote $S_t$ via LLM condicionado em $H$
  \STATE $\forall x \in S_t$: $score(x) = w_1 err(x) + w_2 unc(x) + w_3 div(x)$
  \STATE $H_t \leftarrow \text{top-}k(S_t, score)$
  \STATE $H \leftarrow H \cup H_t$
  \STATE Retreina $f$ em Data 0 $\cup H$
\ENDFOR
\RETURN avalia $f$ no holdout de teste
\end{algorithmic}
\end{figure}

\section{Experimentos}
\subsection{Configura\c{c}\~ao}
Tr\^es bra\c{c}os s\~ao comparados: \textbf{Real} (treina apenas no treino escasso), \textbf{Static} (gera\c{c}\~ao est\'atica: o LLM gera $N$ amostras com prompt gen\'erico e todas s\~ao adicionadas) e \textbf{ActiveInheritance} (m\'etodo proposto, $T{=}5$, lote $64$, top-$k$ $16$). M\'etricas: F1, precis\~ao e revoca\c{c}\~ao no holdout balanceado; diversidade (similaridade de cosseno m\'edia do pool); fidelidade (dist\^ancia de Wasserstein entre embeddings reais e sint\'eticos, projetados nas top-3 componentes PCA dos reais); e custo (tokens acumulados). O Static varia o or\c{c}amento: $N \in \{32, 79, 158\}$.

\subsection{Resultados}
A Tabela~\ref{tab:res} resume o desempenho (m\'edia $\pm$ desvio entre sementes).

\begin{table}[t]
\caption{F1 no teste e custo por bra\c{c}o (m\'edia $\pm$ desvio entre 3 sementes).}
\label{tab:res}
\begin{tabular}{lcc}
\toprule
Bra\c{c}o & F1 & tokens/run \\
\midrule
Real & $0.887 \pm 0.026$ & 0 \\
Static ($N{=}32$) & $0.789 \pm 0.113$ & $\sim$20k \\
Static ($N{=}79$) & $0.852 \pm 0.034$ & $\sim$53k \\
Static ($N{=}158$) & $0.916 \pm 0.013$ & $\sim$103k \\
\textbf{ActiveInheritance} & $\mathbf{0.924 \pm 0.007}$ & $\sim$645k \\
\midrule
Abla\c{c}\~ao pesos (E4) & $0.930$ & $\sim$651k \\
Abla\c{c}\~ao top-$k{=}8$ (E5) & $0.919$ & $\sim$657k \\
\bottomrule
\end{tabular}
\end{table}

\section{Figuras e An\'alise}
A Figura~\ref{fig:f1tokens} compara o F1 em fun\c{c}\~ao dos tokens acumulados: o m\'etodo proposto atinge o maior F1, mas \`a custa de um or\c{c}amento de tokens cerca de $6\times$ maior que o melhor Static. A Figura~\ref{fig:f1cycle} mostra o F1 por ciclo, evidenciando o ganho incremental do pool herdado. A Figura~\ref{fig:wass} mostra que a fidelidade (Wasserstein) \'e semelhante entre os m\'etodos. A Figura~\ref{fig:div} mostra que a diversidade do pool cai ao longo dos ciclos --- o canal de diversidade n\~ao impede o top-$k$ de concentrar amostras similares. A Figura~\ref{fig:box} evidencia a menor vari\^ancia entre sementes do m\'etodo proposto. A Figura~\ref{fig:score} decomp\~oe o escore por ciclo.

\begin{figure}[t]
\centering
\includegraphics[width=\columnwidth]{fig_f1_vs_tokens.png}
\caption{F1 $\times$ tokens acumulados (log). AI atinge o topo com mais tokens.}
\label{fig:f1tokens}
\end{figure}

\begin{figure}[t]
\centering
\includegraphics[width=\columnwidth]{fig_f1_per_cycle.png}
\caption{F1 por ciclo do ActiveInheritance.}
\label{fig:f1cycle}
\end{figure}

\begin{figure}[t]
\centering
\includegraphics[width=\columnwidth]{fig_wasserstein.png}
\caption{Wasserstein real $\times$ sint\'etico (fidelidade).}
\label{fig:wass}
\end{figure}

\begin{figure}[t]
\centering
\includegraphics[width=\columnwidth]{fig_diversity_cycle.png}
\caption{Diversidade do pool (similaridade m\'edia) por ciclo.}
\label{fig:div}
\end{figure}

\begin{figure}[t]
\centering
\includegraphics[width=\columnwidth]{fig_f1_boxplot.png}
\caption{Consist\^encia entre sementes (boxplot de F1).}
\label{fig:box}
\end{figure}

\begin{figure}[t]
\centering
\includegraphics[width=\columnwidth]{fig_score_decomp.png}
\caption{Decomposi\c{c}\~ao do escore (erro, incerteza, diversidade) por ciclo.}
\label{fig:score}
\end{figure}

\section{Discuss\~oes}
Tr\^es achados principais emergem. \textbf{(1) Consist\^encia:} o m\'etodo proposto tem o menor desvio entre sementes ($0.007$), enquanto o Static com $N{=}32$ colapsa em uma semente (F1 $0.63$) --- evid\^encia de que a gera\c{c}\~ao est\'atica em baixo volume pode envenenar a distribui\c{c}\~ao, como alertado por~\cite{souza2025}. \textbf{(2) F1:} o m\'etodo atinge $0.924$, acima do melhor Static ($0.916$) e do Real ($0.887$), embora a diferen\c{c}a n\~ao seja estatisticamente significativa com $n{=}3$ sementes (Wilcoxon $p{=}0.25$; Kruskal-Wallis $p{=}0.15$). \textbf{(3) Efici\^encia:} a hip\'otese de custo n\~ao foi confirmada --- o m\'etodo gera $320$ amostras por run ($\sim$645k tokens) contra $158$ do melhor Static ($\sim$103k tokens) para F1 compar\'avel. Para testar a tese de efici\^encia, seria necess\'ario reduzir o lote do m\'etodo (ex.: lote $32$) para igualar o or\c{c}amento de tokens do Static.

\section{Conclus\~ao}
A Heran\c{c}a Ativa produziu o classificador mais consistente e com o maior F1 no cen\'ario de escassez, mas n\~ao demonstrou vantagem de efici\^encia de tokens na configura\c{c}\~ao atual. Como trabalho futuro, planeja-se: (i) igualar os or\c{c}amentos de tokens entre os bra\c{c}os; (ii) calibrar o canal de diversidade para evitar a concentra\c{c}\~ao do pool; e (iii) avaliar o protocolo TSTR (treinar apenas no sint\'etico).

\footnotesize
\bibliographystyle{IEEEtran}
\begin{thebibliography}{00}

\bibitem{souza2025}
M.~D. de Souza and C.~Zanchettin,
``Breaking the Barrier of Hard Samples: A Data-Centric Approach to Synthetic Data for Medical Tasks,''
in \emph{Proc. ICML}, 2025.

\bibitem{seedat2022}
N.~Seedat, J.~Crabb\'e, and M.~van der Schaar,
``Data-IQ: Characterizing subgroups with heterogeneous outcomes in tabular data,''
in \emph{Proc. NeurIPS}, 2022.

\bibitem{monteiro2018}
R.~Monteiro, R.~Santos, T.~Pardo, T.~de Almeida, E.~Ruiz, and O.~Vale,
``Contributions to the study of fake news in Portuguese: New corpus and automatic detection results,''
in \emph{Proc. PROPOR}, 2018.

\bibitem{souza2020}
F.~Souza, R.~Nogueira, and R.~Lotufo,
``BERTimbau: Pretrained BERT models for Brazilian Portuguese,''
in \emph{Proc. BRACIS}, 2020.

\bibitem{reimers2019}
N.~Reimers and I.~Gurevych,
``Sentence-BERT: Sentence embeddings using Siamese BERT-networks,''
in \emph{Proc. EMNLP}, 2019.

\bibitem{shumailov2024}
I.~Shumailov et al.,
``AI models collapse when trained on recursively generated data,''
\emph{Nature}, vol.~631, 2024.

\bibitem{alaa2022}
A.~Alaa, B.~van Breugel, E.~Saveliev, and M.~van der Schaar,
``How faithful is your synthetic data? Sample-level metrics for evaluating and auditing generative models,''
in \emph{Proc. ICML}, 2022.

\end{thebibliography}

\end{document}
